\documentclass[../main.tex]{subfiles}

\begin{document}
\section{\textbf{Generalized Lagrange Equation}}
\makebox[\textwidth][s]{Author: Xianchao \hfill Date: \today}
Starting from the principle of the least action, which has the definition:
\begin{equation}
  S[t] = \int_{t_1}^{t_2} L(\bm{x}, \dot{\bm{x}}, t)\,\dd{t}
\end{equation}
Here the t is the parameter, usually we use it as time. $\bm{x}$ is the generalized coordination, and $L$ is the Lagrange. $\bm{x}$ is $t$ dependent
which means $\bm{x}(t)$. According to the 
principle of the least action, we have:
\begin{equation}
  \delta S = 0
\end{equation}
We only assume the variation of $x$, which gives:
\begin{align}
  \delta S &= \int_{t_1}^{t_2} \delta L(\bm{x}, \dot{\bm{x}}, t)\,\dd{t}\nonumber\\
  &= \int_{t_1}^{t_2} \frac{\partial L}{\partial x^i}\delta x^i + \frac{\partial L}{\partial \dot{x}^i}\delta\dot{x}^i + \frac{\partial L}{\partial t}\delta t\,\dd{t}\nonumber\\
  &= \int_{t_1}^{t_2} \frac{\partial L}{\partial x^i}\delta x^i + \frac{\partial L}{\partial \dot{x}^i}\delta\dot{x}^i\,\dd{t}
\end{align}
And we have:
\begin{align}
  \int_{t_1}^{t_2} \frac{\partial L}{\partial\dot{x}^i}\delta\dot{x}^i \,\dd{t} &=  \left.\left[\frac{\partial L}{\partial \dot{x}^i}\delta x^i\right]\right|^{t_2}_{t_1}
  - \int_{t_1}^{t_2}\dv{t}(\frac{\partial L}{\partial \dot{x}^i})\delta x^i\,\dd{t}  \nonumber\\
  &= -\int_{t_1}^{t_2}\dv{t}(\frac{\partial L}{\partial \dot{x}^i})\delta x^i\,\dd{t} 
\end{align}
And now we have:
\begin{equation}
  \delta S = \int_{t_1}^{t_2} \frac{\partial L}{\partial x^i}\delta x^i - \dv{t}(\frac{\partial L}{\partial \dot{x}^i})\delta x^i\,\dd{t} = 0
\end{equation}
Thus we have the Euler-Lagrange equation:
\begin{equation}
  \dv{t}(\frac{\partial L}{\partial \dot{x}^i}) - \frac{\partial L}{\partial x^i} = 0
\end{equation}
We can define $\frac{\partial L}{\partial \dot{x}^i} = p^i$, applying Legendre transformation, we have Hamiltonia:
\begin{equation}
  H(\bm{x}, \bm{p}) = p_i\dot{x}^i - L(\bm{x}, \bm{p})
\end{equation}
Actually the term $p_i\dot{x}^i$ is the so called symplectic potential, and we can replace the $\bm{x}$ and $\bm{p}$ 
in a very general form $\bm{\xi}$, so the Lagrange can be written as:
\begin{equation}
  L = -a_i(\bm{\xi})\dot{\xi}^i - H(\bm{\xi})
\end{equation}
Using this form, we can substitute it into the Euler-Lagrange equation, which has:
\begin{align}
  \pdv{L}{\xi^i} &= -\pdv{a_j}{\xi^i}\dot{\xi}^j - \pdv{H}{\xi^i}\\
  \pdv{L}{\dot{\xi}^i} &= -a_i 
\end{align}
So we have:
\begin{align}
  \dv{t}(\frac{\partial L}{\partial \dot{\xi}^i}) - \frac{\partial L}{\partial \xi^i} &= -\pdv{a_i}{\xi^j}\dot{\xi}^j + \pdv{a_j}{\xi^i}\dot{\xi}^j + \pdv{H}{\xi^i}\nonumber\\
  &= \left(\pdv{a_j}{\xi^i} - \pdv{a_i}{\xi^j}\right)\dot{\xi}^j + \pdv{H}{\xi^i} = 0
\end{align}
Here we define the term $\pdv{a_j}{\xi^i} - \pdv{a_i}{\xi^j}$ as $b_{ij}$, then we have:
\begin{equation}
  b_{ij}\dot{\xi}^j = -\pdv{H}{\xi^i}
\end{equation}
Here the $\bm{a}(\bm{\xi})$ is the Berry connection, while the $b_{ij}$ is the Berry curvature.

For quantum system, here we only consider the coherent state, which means that the quantum state can be uniquely
determined by a set of parameters, which is $\ket{\psi(\bm{\xi})}$. In this case the single point in phase space is
exactly a quantum state.
The Schrodinger equation is:
\begin{equation}
  i\dv{t}\ket{\psi} = \hat{H}\ket{\psi}
\end{equation}
Which is obtained from the definition of action:
\begin{equation}
  S = \int_{t_1}^{t_2} \bra{\psi} i\dv{t} - \hat{H} \ket{\psi}
\end{equation}
And we can get:
\begin{equation}
  i\bra{\psi}\dv{t}\ket{\psi} - \bra{\psi} \hat{H}\ket{\psi} = 0
\end{equation}
Comparing with the generalized Lagrange equation, we have:
\begin{equation}
  i\bra{\psi}\dv{t}\ket{\psi} = -a_i(\xi)\dot{\xi}^i
\end{equation}
Now look into the term $i\bra{\psi}\dv{t}\ket{\psi}$, which is:
\begin{equation}
  i\bra{\psi}\dv{t}\ket{\psi} = i\bra{\psi}\pdv{\xi^i}\dot{\xi}^i\ket{\psi}
\end{equation}
$\dot{\xi}^i$ and $\ket{\psi}$ are commutable, so we have:
\begin{equation}
  i\bra{\psi}\dv{t}\ket{\psi} = i\bra{\psi}\pdv{\xi^i}\ket{\psi}\dot{\xi}^i
\end{equation}
Thus we have:
\begin{equation}
  a_i = -i\bra{\psi}\pdv{\xi^i}\ket{\psi}
\end{equation}
\end{document}
